% =============================================================
% บทที่ 5 - สรุปผลและอภิปรายผล
% โครงงาน: การออกแบบและพัฒนาเกมเดินเรือด้วยระบบสั่งการด้วยเสียงเพื่อส่งเสริมทักษะการสื่อสาร
% =============================================================

\section{สรุปผลการดำเนินงาน}

รายงานความก้าวหน้าฉบับนี้สรุปผลการดำเนินงานในระยะที่ 1 ของโครงงานการออกแบบและพัฒนาเกมเดินเรือด้วยระบบสั่งการด้วยเสียง (Voice-Controlled Ship Navigation Game) เพื่อส่งเสริมทักษะการสื่อสาร ซึ่งครอบคลุมการศึกษาทฤษฎีและงานวิจัยที่เกี่ยวข้อง การออกแบบระบบและกลไกการเล่น และการพัฒนาเกมต้นแบบ (Prototype) โดยผู้จัดทำได้ดำเนินการแล้วเสร็จในส่วนต่าง ๆ ดังนี้

ด้านการออกแบบ ผู้จัดทำได้ออกแบบเนื้อเรื่องและสตอรี่บอร์ด กระบวนการเล่นหลัก (Core Game Loop) สถาปัตยกรรมของระบบ ส่วนติดต่อผู้ใช้ (UI/HUD) และบทบาทของตัวละครทั้งสามตำแหน่ง ได้แก่ กัปตัน (ผู้เล่น) พลต้นหน และพลถือท้าย รวมถึงจัดทำคลังคำสั่งการเดินเรือมาตรฐาน (Standard Marine Communication Phrases) สำหรับใช้เป็นเกณฑ์ตรวจสอบความถูกต้องของคำสั่งเสียง ดังแสดงในภาคผนวก ง

ด้านการพัฒนา ได้พัฒนาเกมต้นแบบบนเอนจิน Unity 3D ด้วยภาษา C\# และเชื่อมต่อระบบย่อยหลักได้สำเร็จ 3 ระบบ ได้แก่ (1) ระบบรับคำสั่งเสียง ซึ่งรับสัญญาณจากไมโครโฟนและแปลงเป็นข้อความผ่าน Whisper API (2) ระบบลูกเรือปัญญาประดิษฐ์ ซึ่งส่งข้อความคำสั่งไปตีความผ่าน OpenAI API และสร้างบทสนทนาโต้ตอบกลับ และ (3) ระบบฟิสิกส์และการควบคุมเรือ ซึ่งจำลองความเฉื่อย (Inertia) และรัศมีวงเลี้ยว (Turning Radius) ของเรือขนาดใหญ่ ผลการทดสอบเบื้องต้นในระดับหน่วย (Unit Testing) พบว่าระบบสามารถรับคำสั่งเสียง แปลงเป็นข้อความ ส่งต่อไปยังตัวแยกคำสั่ง (Command Router) และควบคุมการเคลื่อนที่ของเรือได้อย่างต่อเนื่องตามที่ออกแบบไว้ รวมถึงสามารถจำลองสถานการณ์สื่อสารติดขัด (Communication Breakdown) โดยให้ลูกเรือ AI ทวนถามเมื่อคำสั่งไม่ชัดเจนได้ตามกลไกที่กำหนด

งานที่อยู่ระหว่างดำเนินการ ได้แก่ การพัฒนาระบบสุ่มสร้างสิ่งกีดขวาง (Obstacle Generation System) และระบบอินเทอร์เฟซผู้ใช้ (UI/HUD) ให้สมบูรณ์ การปรับจูนระดับความยากและค่าความคลาดเคลื่อนที่ยอมรับได้ (Tolerance) ของการรับคำสั่งเสียง ตลอดจนการทดสอบแบบบูรณาการ (Integration Testing) และการทดสอบการยอมรับของผู้ใช้ (User Acceptance Testing)

ในส่วนของการประเมินผล ยังไม่ได้ดำเนินการเก็บข้อมูลในระยะนี้ โดยผู้จัดทำมีแผนจะประเมินประสิทธิภาพทางเทคนิค (ความแม่นยำและความหน่วงของระบบ Speech-to-Text) และประเมินประสบการณ์ผู้ใช้ ความมั่นใจ และระดับความวิตกกังวลในการใช้ภาษาอังกฤษ กับกลุ่มตัวอย่างจำนวน 40–50 คน ในระยะถัดไป ซึ่งผลการประเมินดังกล่าวจะเป็นข้อมูลสำหรับสรุปว่าระบบบรรลุวัตถุประสงค์ในการลดความวิตกกังวลในการใช้ภาษาต่างประเทศ (Foreign Language Anxiety) ได้มากน้อยเพียงใด

\section{ปัญหาและข้อจำกัด}

ในการดำเนินโครงงานและพัฒนาตัวเกมต้นแบบ ผู้จัดทำพบปัญหาและข้อจำกัดที่ส่งผลต่อประสิทธิภาพของระบบและการประเมินผล ดังนี้

\begin{enumerate}[label=\arabic*)]
    \item \textbf{ข้อจำกัดด้านฮาร์ดแวร์และสภาพแวดล้อม:} ประสิทธิภาพและความแม่นยำในการรับคำสั่งของเกมขึ้นอยู่กับคุณภาพของไมโครโฟนที่ผู้เล่นใช้งาน รวมถึงหากสภาพแวดล้อมขณะเล่นมีเสียงรบกวนภายนอก (Background Noise) มากเกินไป จะส่งผลให้ระบบรับคำสั่งผิดพลาดได้ง่าย
    \item \textbf{ข้อจำกัดด้านระบบการรู้จำเสียง:} ระบบ Speech-to-Text มีข้อจำกัดในการทำความเข้าใจสำเนียง (Accent) ที่หลากหลายของผู้เล่นที่ไม่ใช่เจ้าของภาษา (Non-native speakers) 
    \item \textbf{ข้อจำกัดด้านขอบเขตการสื่อสาร:} โครงงานนี้ออกแบบมารองรับเฉพาะชุดคำสั่งภาษาอังกฤษเกี่ยวกับการเดินเรือที่กำหนดไว้ล่วงหน้าเท่านั้น ไม่รองรับการสนทนาโต้ตอบแบบอิสระ (Free-form Conversation)
    \item \textbf{ข้อจำกัดด้านกลุ่มตัวอย่างและการประเมินผล:} การประเมินผลอาศัยกลุ่มตัวอย่างจำนวนน้อย (40–50 คน) ภายในบริบทของมหาวิทยาลัย ซึ่งอาจไม่สะท้อนผลลัพธ์ของประชากรในวงกว้างทั้งหมดได้ และเกมเป็นเพียงเครื่องมือเสริม ไม่สามารถใช้ทดแทนการทดสอบมาตรฐานทางภาษาเชิงลึกได้
\end{enumerate}

\section{แนวทางการแก้ไขและข้อเสนอแนะ}

เพื่อให้การพัฒนาเกมมีความสมบูรณ์และสามารถนำไปต่อยอดได้ในอนาคต ผู้จัดทำมีแนวทางการแก้ไขและข้อเสนอแนะ ดังนี้

\begin{enumerate}[label=\arabic*)]
    \item \textbf{การจัดการเสียงรบกวน:} ผู้พัฒนาควรแนะนำให้ผู้เล่นใช้งานชุดหูฟังที่มีไมโครโฟนในตัว หรืออาจต้องศึกษาการพัฒนาระบบกรองเสียงรบกวนเบื้องต้น (Noise Filter) ก่อนส่งข้อมูลเสียงไปยัง API เพื่อเพิ่มความแม่นยำในการรับคำสั่ง
    \item \textbf{การปรับปรุงความแม่นยำในการรับคำสั่ง:} ควรมีการรวบรวมข้อมูล หรือปรับค่าความคลาดเคลื่อน (Tolerance) ของคำสั่งที่อนุโลมได้ เพื่อรองรับสำเนียงที่หลากหลายและลดความหงุดหงิดของผู้เล่น
    \item \textbf{การขยายขอบเขตในอนาคต:} ในอนาคตสามารถขยายระบบคลังคำศัพท์ (Command Vocabulary) ให้รองรับรูปแบบการสื่อสารทางทะเลที่ซับซ้อนยิ่งขึ้น และเพิ่มการประเมินผลกับกลุ่มตัวอย่างที่ใหญ่ขึ้น เพื่อประเมินผลสัมฤทธิ์ด้านการลดตัวกรองทางอารมณ์ (Affective Filter) ในระดับที่ครอบคลุมกว่าเดิม
\end{enumerate}